\cotefile{134-2}
\scanpage{178}\mspage{33}
vérifier et d'expliciter comment les ``6 opérations'' habituelles sur les faisceaux (soit sur $X$, soit sur une réunion de strates de $X$) peuvent s'exprimer dans ce dictionnaire, dans le cas disons de strates non singulières (\unsure{autrement}, on aurait des difficultés avec les complexes dualisants, dont on aimerait qu'ils appartiennent à notre catégorie\unsure{.\,.\,.\,.\,.|\ldots}) \inserted{et de revérifier le formulaire connu de ces opérations}. Mais il \unsure{apparaitra|apparaîtra|apparaît} probablement que, pour bien pouvoir tout transcrire, il convient, plutôt \inserted{que} de regarder un diagramme du type
\diagram{\begin{tikzcd}[ampersand replacement=\&,column sep=1.4em,row sep=1.2em]
\cdot \arrow[r] \arrow[d] \& {} \\
{} \& {}
\end{tikzcd}}
dans la catégorie homotopique formée avec la cat. des ens. semi-simpliciaux, \inserted{de} regarder la catégorie des diagrammes d'ensembles $1/2$ simpliciaux, et à partir de là passer à une catégorie homotopique de fractions.

J'ai plus ou moins explicité dernièrement, en cogitant sur \deleted{\unsure{les|des|ces}}\unsure{des|les} fondements de ``topologie modérée'' (i.e. où on élimine d'emblée tous les phénomènes de sauvagerie), comment une stratification équisingulière\inserted{\deleted{(\unsure{d'un} \unsure{espace} \unsure{modéré} \unsure{compact})}}, à \unsure{strates} non singulières disons, d'un ``espace modéré'' \inserted{compact}, donne naissance canoniquement à un diagramme \deleted{\unsure{fon|f}} d'espaces,
% archive sheet number: 178
\scanpage{179}\mspage{34}
qui sont des variétés à bord, les \deleted{\unsure{morp|morph}} flèches du diagramme étant pour l'essentiel des fibrations loc. triviales de variétés à bord sur d'autres (à fibres des variétés à bord), \uline{et} \deleted{\unsure{les|des}}\unsure{des|les} inclusions des bords dans les variétés à bord \deleted{(}[en fait, on trouve des inclusions légèrement plus générales, certains des bords qui apparaissent étant munis d'une décomposition cellulaire ``élémentaire'', \deleted{\unsure{en cellules|en éléments}}\inserted{i.e. \unsure{dont}} \unsure{toutes} les strates \inserted{fermées} sont encore des variétés à bord qui se recollent suivant des parties communes du bord~; et il y a lieu d'envisager \inserted{aussi} les inclusions de ces morceaux les uns dans les autres\ldots.\deleted{)}], et se réconstitue à partir de ce diagramme \inserted{par recollement}. En d'autres termes, on a une description-dévissage canonique des espaces modérés compacts \inserted{$X$}, éventuellement munis de stratifications équisingulières à strates non sing., en termes de diagrammes \inserted{finis} de nature précisée formés avec des variétés à bord. Quand on s'intéresse aux faisceaux \inserted{(ou complexes de faisceaux, ou $n$-champs\ldots)} $\unsure{P|\mathcal{P}}$-constructibles sur $X$, où $P$ est une telle stratification fixée, ceux-ci s'explicitent en termes du diagramme envisagé, dont il
% archive sheet number: 179
\scanpage{180}\mspage{35}
n'y a plus lieu de retenir que le ``type d'homotopie''. On prévoit que les 6 opérations sur ces faisceaux peuvent se traduire ad-hoc dans ce cadre homotopique. Enfin, si au lieu d'un seul espace modéré compact $X$, on a disons un morphisme modéré $f:X\unsure{-|\to}Y$ de tels, alors, choisissant des stratifications équisingulières sur $X$ et $Y$ adaptées à $f$ (les strates de $X$ étant notamment des fibrés loc. triviaux sur celles de $Y\ldots$), on doit trouver un ``morphisme'' du diagramme de variétés à bord exprimant $X$ dans celui exprimant $Y$ (avec des morphismes verticaux se ramenant \deleted{\unsure{à|,}} encore, pour l'essentiel, à des fibrations de variétés à bord sur d'autres) de telle façon que les 4 opérations $\mathrm{R}\deleted{\unsure{f}}g_*$, $\mathrm{R}g_!$, $\mathrm{L}g^*$, $g^!$ entre faisceaux \inserted{$\unsure{P|\mathcal{P}}_X$- et $P_Y$-constructibles} sur $X$ et $Y$ (ou sur des sous-espaces loc. fermés \inserted{$X'$, $Y'$} réunions de strates, de sorte que $f$\unsure{,|.} induise $g:X'\unsure{-|\to}Y'$) s'expriment encore en termes d'opérations standart entre \unsure{seuls|de tels} types d'homotopie. Enfin, toutes ces constructions, encore partiellement hypothétiques (y a du travail \inserted{de fondements} sur la planche!) doivent
% archive sheet number: 180
\scanpage{181}\mspage{36}
pouvoir se paraphraser dans le cadre des schémas excellents, à grands coups de voisinages tubulaires étales. Dans l'un et l'autre cas, le ``type d'homotopie fin'' d'un espace modéré, resp. d'un schéma excellent, se définirait par passage à la limite à partir des ``\inserted{$\unsure{P|\mathcal{P}}$-}types d'homotopie'' associés à des stratifications équisingulières (à strates non sing.) $P$ de plus en plus fines. Ce type d'homotopie \inserted{``fin''} incarnerait la connaissance, non des seuls faisceaux ou \inserted{$n$-}champs loc constants, mais \inserted{(via un passage à la limite \unsure{inductive})} la connaissance de \uline{tous}. Et il dépendrait, en un sens convenable, fonctoriellement de $X$. Dans le cas d'un schéma de type fini sur un corps alg. clos disons, les \uline{théorèmes de finitude} cohomologique ou homotopique les plus forts s'exprimeraient précisément en termes des types d'homotopie fin, et diraient que les types d'homotopie \unsure{ordinaire} qui en sont les \unsure{élément|éléments} constituants sont essentiellement des ``\unsure{polyédres|polyèdres} finis'' -- et \unsure{$\overline{m}$|$=$} des variétés à bord compactes -- ou de façon plus précise,
% archive sheet number: 181
\scanpage{182}\mspage{37}
les ``complétés profinis'' \inserted{(au sens d'Artin-Mazur)} premiers à $p$ (la car. de $k$) de tels \unsure{espaces}. On voit clairement comment aborder un tel programme en car. $0$, mais on prévoit de l'amusement supplémentaire, voire du mystère, en car $p>0$, pour les variétés qui, même birationnellement, refuseraient de se laisser remonter en car. $0$!

\begin{center}\rule{0.18\linewidth}{0.4pt}\end{center}

De ces cogitations essentiellement géométriques je ne saurais à l'heure actuelle tirer un programme précis de développement de structures algébriques adéquates pour les exprimer. Je me bornerai à quelques remarques en marge.

\deleted{\unsure{a)|$\alpha$)}} Depuis longtemps j'ai été \unsure{intrigu-|intrigué} par l'idée d'une ``linéarisation'' d'un type d'homotopie (ordinaire), i.e. la question du type: Si $X$ est un type d'homotopie, combien d'information cohomologique du type: cohomologie de $X$ à coefficients variables $M$ (\unsure{constant|constants} ou constant tordus), structure multiplicative $\deleted{\unsure{+|=}}H^i(X,M)\times H^j(X,N)\unsure{-|\to}H^{i\unsure{\cdot|+}j}(X,M\otimes N)$, plus éventuellement d'autres opérations cohomologiques -- est-il nécessaire d'avoir pour réconstituer le type d'homotopie entièrement \deleted{\unsure{?}}? (disons,
% archive sheet number: 182
\scanpage{183}\mspage{38}
dans cette première approche pré-catégorie dérivée, par la connaissance du groupe fondamental $\pi_1$, \unsure{donc} des catégories des coeff. \unsure{constant|constants} tordus $=$ $\pi_1$-\unsure{modules|module}, les foncteurs $H^i(-,M)$ sur ceux-ci avec la structure de foncteur \deleted{\unsure{ab|ab.}, de \unsure{foncteur}} \inserted{cohomologique} relatif aux suites exactes, la \deleted{\unsure{struc}}structure de cup-produit etc -- reliés par certaines propriétés formelles. Une fois qu'on dispose du langage des cat. dérivées, la \inserted{sous catégorie de la catégorie dérivée des} complexes abéliens, \unsure{form-|formée|formés} \inserted{aux complexes de} faisceaux de coh. loc constante sur $X$, avec sa structure \unsure{triangulé|triangulée} et sa structure \deleted{\unsure{de|sou|m}}multiplicative $\overset{\mathbb{L}}{\otimes}$ (et éventuellement $\mathbb{R}\mathrm{Hom}$ \ldots) donnait un candidat plus satisfaisant pour espérer récupérer le type d'homotopie. Je n'ai guère d'idée si oui ou non ils permettent de le récupérer \inserted{en effet}, mais par contre n'ai guère de doute qu'en allant au bout de la ``linéarisation'', c'est à dire en partant dans le \uline{non abélien} et travaillant avec la $(n+1)$-catégorie (\unsure{sans|sa} \unsure{structures|structure} supplémentaires \unsure{dessus|disons}!) des $n$-champs loc. constants \inserted{\unsure{sur $X$}} pour tout $n$
% archive sheet number: 183
\scanpage{184}\mspage{39}
on arrive à \unsure{réconstituer|reconstituer} le type d'homotopie via son $\infty$-groupoïde fondamental, comme expliqué dans ma lettre précédente, \deleted{\unsure{e|a}}et rappelé dans celle-ci. (Cela signifie en particulier que toutes les opérations cohomologiques possibles et imaginables sont déjà incluses dans la donnée d'un tel système de $n$-catégories .\,.\,.\,.) \deleted{Mais}

De même, les types d'homotopie plus \unsure{cannelés|canelés}, liés à certains diagrammes finis, qu'on peut associer à \deleted{\unsure{des}}certains types de stratifications $\unsure{P|\mathcal{P}}$ d'espaces topologiques ``modérés'' $X$ disons, devraient correspondre de façon aussi parfaite aux \deleted{\unsure{$n$-}}\inserted{$(n+1)$-}catégories de $n$-champs sur $X$ qui sont loc constants sur chacune des strates de $P$ (disons: qui sont \unsure{subordonnés} à $P$). Si plus haut le \unsure{topos|topis|topo} de la cat. dérivée loc. constante marchait bel et bien, on s'attendrait ici à récupérer le type d'homotopie mixte à partir de la \inserted{\unsure{ss-}}catégorie de la \unsure{catégori|catégorie} dérivée des faisceaux abéliens sur $X$, formée des complexes qui sont à ch. loc. constante sur chacune des strates -- avec en plus les opérations $\overset{\mathbb{L}}{\otimes}$, $\mathbb{R}\underline{\mathrm{Hom}}$ plus, les
% archive sheet number: 184
\scanpage{185}\mspage{40}
\unsure{cas|Cas} échéant, les 4 opérations $\mathbb{R}g_!$ $\mathbb{R}g_*$, $\mathrm{L}g^*$, $g^!$ pour des inclusions $g:Z'-Z''$ de réunions localement fermées de strates\unsure{.\,.\,.\,.\,.\,.|.\,.\,.\,.\,.} (L'ennui ici est qu'on ne sait même pas encore ce que c'est qu'une catégorie triangulée, pas plus qu'on ne sait ce qu'est la version non commutative, sans doute plus simple et plus fondamentale: \uline{une} ``catégorie homotopique'', avec opérations ``fibre'' et ``cofibre'' .\,.\,.\,.)

\begin{center}\rule{0.41\linewidth}{0.4pt}\end{center}

Il est vraiment temps que je termine cette lettre-fleuve, qui devient de plus en plus vaseuse. Juste une question: Quelle est cette mirifique formule de \unsure{Bloch|Blod}-Quillen à laquelle tu fais allusion, dont je n'ai jamais entendu parler, et qui fait venir l'eau à la bouche?

Bien cordialement à toi\\
\hspace*{0.42\linewidth}Alexandre\\
\hspace*{0.52\linewidth}Grothendieck\\
% archive sheet number: 185
